\documentclass[11pt]{article}
\usepackage{mathblatt}
% gesetzt von werkzeuge/setzer.py; Lösungen nach bau/bauregeln.md „Lösungen“.
\newcommand{\kennung}{HMF}
\begin{document}
\pfheftstil
\fancyfoot[C]{{\footnotesize\color{mbgrau}\kennung\hfill\thepage}}
\pfheftkopf{Vergleichen, Ordnen, Runden}{Klasse 6 \textperiodcentered{} Lösungen}

\begin{pfloesung}{}
\lz{1b)}{a) $>$ \quad b) $>$ \quad c) $<$ \quad d) $>$}{a) $0{,}58 > 0{,}54$; b) $4{,}62 > 4{,}26$; c) $15{,}2 < 16{,}2$; d) $0{,}071 > 0{,}017$}{}
\end{pfloesung}
\begin{pfloesung}{}
\lz{2.}{a) $>$ \quad b) $<$ \quad c) $>$ \quad d) $=$}{a) $0{,}70 > 0{,}69$; b) $3{,}45 < 3{,}50$; c) $0{,}410 > 0{,}405$; d) $5{,}10 = 5{,}10$}{}
\end{pfloesung}
\begin{pfloesung}{}
\lz{3.}{$4{,}03 < 4{,}3 < 4{,}303 < 4{,}33$}{$4{,}03$; $4{,}3$; $4{,}303$; $4{,}33$ – mit Nullen: $4{,}030 < 4{,}300 < 4{,}303 < 4{,}330$}{}
\end{pfloesung}
\begin{pfloesung}{}
\lz{4.}{$5{,}8 > 5{,}088 > 5{,}08 > 5{,}008$}{$5{,}8$; $5{,}088$; $5{,}08$; $5{,}008$ – mit Nullen: $5{,}800 > 5{,}088 > 5{,}080 > 5{,}008$}{}
\end{pfloesung}
\begin{pfloesung}{}
\lz{5.}{Silber Ida ($41{,}1$\,s), Bronze Gül ($41{,}75$\,s)}{Silber: Ida ($41{,}1$\,s), Bronze: Gül; die kürzeste Zeit gewinnt: $41{,}08 < 41{,}10 < 41{,}75 < 41{,}80 < 42{,}10$ (Finn, Ida, Gül, Ela, Hannes).}{}
\end{pfloesung}
\begin{pfloesung}{}
\lz{6b)}{a) $2{,}7$ \quad b) $5{,}5$ \quad c) $0{,}4$ \quad d) $8{,}0$}{a) $2{,}73 \approx 2{,}7$; b) $5{,}48 \approx 5{,}5$; c) $0{,}35 \approx 0{,}4$; d) $8{,}04 \approx 8{,}0$}{}
\end{pfloesung}
\begin{pfloesung}{}
\lz{7.}{a) $1{,}23$ \quad b) $6{,}08$ \quad c) $12{,}47$ \quad d) $0{,}56$}{a) $1{,}234 \approx 1{,}23$; b) $6{,}078 \approx 6{,}08$; c) $12{,}465 \approx 12{,}47$; d) $0{,}555 \approx 0{,}56$}{}
\end{pfloesung}
\begin{pfloesung}{}
\lz{8.}{a) $3{,}0$ \quad b) $10{,}0$ \quad c) $4{,}00$ \quad d) $1{,}0$}{a) $2{,}97 \approx 3{,}0$; b) $9{,}96 \approx 10{,}0$ (9 Zehntel $+ 1 =$ 1 Einer; 9 Einer $+ 1 = 10$); c) $3{,}998 \approx 4{,}00$; d) $0{,}95 \approx 1{,}0$}{}
\end{pfloesung}
\begin{pfloesung}{}
\lz{9.}{a) $2{,}35$\,€ \quad b) $4{,}0$\,kg \quad c) $1{,}78$\,m \quad d) $12{,}1$\,s}{a) $2{,}346$\,€ $\approx 2{,}35$\,€; b) $3{,}96$\,kg $\approx 4{,}0$\,kg; c) $1{,}784$\,m $\approx 1{,}78$\,m; d) $12{,}05$\,s $\approx 12{,}1$\,s}{}
\end{pfloesung}
\begin{pfloesung}{}
\lz{10.}{Nein: Der Bon zeigt $4{,}87$\,€.}{Nein; auf Cent (Hundertstel) runden: rechts steht die 5, also aufrunden: $4{,}865$\,€ $\approx 4{,}87$\,€; $4{,}87$\,€ $> 4{,}85$\,€, es fehlen 2 Cent.}{}
\end{pfloesung}
\begin{pfloesung}{}
\lz{11b)}{a) z.\,B. $0{,}35$ \quad b) z.\,B. $3{,}45$ \quad c) z.\,B. $0{,}255$}{a) z.\,B. $0{,}35$; b) z.\,B. $3{,}45$; c) z.\,B. $0{,}255$}{}
\end{pfloesung}
\begin{pfloesung}{}
\lz{12.}{a) $0{,}65$ \quad b) $4{,}15$ \quad c) $0{,}045$}{a) $60$ und $70$ $\to$ $65$, Mitte $= 0{,}65$; b) $410$ und $420$ $\to$ $415$, Mitte $= 4{,}15$; c) $40$ und $50$ (Tausendstel) $\to$ $45$, Mitte $= 0{,}045$}{}
\end{pfloesung}
\begin{pfloesung}{}
\lz{13.}{a) $1{,}5$ \quad b) $1{,}35$}{a) $120$ und $180$ $\to$ $150$, Mitte $= 1{,}5$ (auf dem Strich); b) $120$ und $150$ $\to$ $135$, Mitte $= 1{,}35$ (zwischen $1{,}3$ und $1{,}4$)}{}
\end{pfloesung}
\begin{pfloesung}{}
\lz{14.}{a) $3{,}05$ \quad b) $6{,}05$}{a) $290$ und $320$, Abstand $30$, $290 + 15 = 305$, Mitte $= 3{,}05$; b) $570$ und $640$, Abstand $70$, $570 + 35 = 605$, Mitte $= 6{,}05$}{}
\end{pfloesung}
\begin{pfloesung}{}
\lz{15.}{Bank bei km $3{,}75$; Rastplatz zwischen Bank und Brücke}{Zwischen Bank und Brücke; Bank: Mitte $= 3{,}75$ ($340$ und $410$, Abstand $70$, $340 + 35 = 375$); Rastplatz km $3{,}80 > 3{,}75$.}{}
\end{pfloesung}
\begin{pfloesung}{}
\lz{16b)}{a) $0{,}6$ \quad b) $0{,}45$ \quad c) $0{,}04$ \quad d) $0{,}07$ \quad e) $0{,}75$ \quad f) $0{,}36$}{a) $\frac{3}{5} = \frac{6}{10} = 0{,}6$; b) $45\,\% = \frac{45}{100} = 0{,}45$; c) $0{,}2 \cdot 0{,}2 = 0{,}04$; d) $7\,\% = \frac{7}{100} = 0{,}07$; e) $\frac{3}{4} = \frac{75}{100} = 0{,}75$; f) $0{,}6 \cdot 0{,}6 = 0{,}36$}{}
\end{pfloesung}
\begin{pfloesung}{}
\lz{17.}{a) $>$ \quad b) $<$ \quad c) $<$ \quad d) $=$}{a) $0{,}75 > 0{,}70$; b) $0{,}08 < 0{,}80$; c) $0{,}49 < 0{,}70$; d) $0{,}15 = 0{,}15$}{}
\end{pfloesung}
\begin{pfloesung}{}
\lz{18.}{$0{,}6^2 < 0{,}38 < 39\,\% < \dfrac{2}{5}$}{$0{,}6^2 = 0{,}36 < 0{,}38 < 39\,\% = 0{,}39 < \frac{2}{5} = 0{,}4$}{}
\end{pfloesung}
\begin{pfloesung}{}
\lz{19.}{$6\,\% = 0{,}06$}{$\frac{6}{100} = 0{,}06$; falsch sind $0{,}3^2 = 0{,}09$, $\frac{1}{4} = 0{,}25 < 0{,}3$, $0{,}90 > 0{,}85$.}{}
\end{pfloesung}
\begin{pfloesung}{}
\lz{20.}{6a ($0{,}75$)}{$\frac{3}{4} = \frac{75}{100} = 0{,}75$; 6b: $0{,}68$; 6c: $72\,\% = 0{,}72$; $0{,}75$ ist am größten.}{}
\end{pfloesung}
\begin{pfloesung}{}
\lz{21.}{$2{,}078 < 2{,}08 < 2{,}18 < 2{,}8$}{$2{,}078$; $2{,}08$; $2{,}18$; $2{,}8$ – mit Nullen: $2{,}078 < 2{,}080 < 2{,}180 < 2{,}800$}{}
\end{pfloesung}
\begin{pfloesung}{}
\lz{22.}{a) $7{,}3$ \quad b) $7{,}35$ \quad c) $10{,}0$}{a) $7{,}348 \approx 7{,}3$; b) $7{,}348 \approx 7{,}35$; c) $9{,}97 \approx 10{,}0$}{}
\end{pfloesung}
\begin{pfloesung}{}
\lz{23.}{$4{,}75$}{$470$ und $480$ $\to$ $475$, Mitte $= 4{,}75$}{}
\end{pfloesung}
\begin{pfloesung}{}
\lz{24.}{a) $=$ \quad b) $<$ \quad c) $<$}{a) $\frac{35}{100} = 0{,}35$, also $=$; b) $0{,}06 < 0{,}60$; c) $0{,}81 < 0{,}90$}{}
\end{pfloesung}
\begin{pfloesung}{}
\lz{25.}{$0{,}2^2 = 0{,}04$}{$0{,}2^2 = 0{,}04$; die anderen sind $0{,}8$, $0{,}08$ und $0{,}07$.}{}
\end{pfloesung}
\begin{pfloesung}{}
\lz{26.}{Lea: Krug B ($0{,}2$); am süßesten Krug A ($0{,}25$)}{Lea nimmt Krug B, am süßesten ist Krug A; A: $\frac{1}{4} = \frac{25}{100} = 0{,}25$; B: $0{,}20$; C: $22\,\% = 0{,}22$; $0{,}20 < 0{,}22 < 0{,}25$.}{}
\end{pfloesung}
\end{document}
